\documentclass[12pt,a4paper]{article}

\usepackage[spanish,es-nodecimaldot]{babel}
\usepackage[utf8]{inputenc}
\usepackage[T1]{fontenc}
\usepackage{lmodern}
\usepackage{geometry}
\geometry{margin=2.5cm}
\usepackage{setspace}
\onehalfspacing
\usepackage{hyperref}
\hypersetup{colorlinks=true,linkcolor=blue,urlcolor=blue}
\usepackage{booktabs}
\usepackage{tabularx}
\usepackage{longtable}
\usepackage{array}
\usepackage{enumitem}
\usepackage{float}
\usepackage{listings}
\usepackage{xcolor}
\usepackage{tikz}
\usetikzlibrary{arrows.meta,positioning,fit,calc,shapes.geometric,shapes.misc,shadows,backgrounds}

\definecolor{boxblue}{RGB}{220,232,248}
\definecolor{boxgreen}{RGB}{214,238,218}
\definecolor{boxorange}{RGB}{253,235,208}
\definecolor{arrowgray}{RGB}{80,80,80}
\definecolor{borderblue}{RGB}{70,130,180}
\definecolor{bordergreen}{RGB}{46,125,50}
\definecolor{borderorange}{RGB}{188,100,0}
\definecolor{lanecolor}{RGB}{245,248,252}
\definecolor{laneheader}{RGB}{60,90,140}

\title{\textbf{Informe de Arquitectura de Microservicios --- Parcial 1}\\[4pt]
{\large Grupo Cordillera: Plataforma de Monitoreo Inteligente}}
\author{Proyecto Semestral}
\date{\today}

\newcolumntype{Y}{>{\raggedright\arraybackslash}X}

\lstdefinestyle{yamlstyle}{
  basicstyle=\ttfamily\small,
  keywordstyle=\color{blue},
  commentstyle=\color{gray},
  frame=single,
  breaklines=true
}

\begin{document}
\maketitle
\thispagestyle{empty}
\clearpage
\tableofcontents
\clearpage

% ─────────────────────────────────────────────────────────────
\section{Resumen Ejecutivo}

Este informe documenta la arquitectura de microservicios propuesta para Grupo Cordillera, en el marco del \textbf{Parcial 1}. El problema concreto es simple de enunciar: la organización consolida datos de ventas, inventario, finanzas y atención al cliente de forma manual, lo que genera demoras, errores y falta de visibilidad para la alta dirección.

La solución se organiza en cuatro servicios backend: \textbf{API Gateway}, \textbf{Auth Service}, \textbf{Data Ingestion Service} y \textbf{KPI Engine}. Los tres primeros están implementados; el cuarto es el siguiente paso. Se aplican los patrones \textit{Repository}, \textit{Factory Method} y \textit{Circuit Breaker}, con orquestación en \texttt{docker-compose}.

% ─────────────────────────────────────────────────────────────
\section{Contexto del Problema}

Grupo Cordillera integra información desde POS, e-commerce, inventario, finanzas y CRM. El proceso es hoy completamente manual y trae consecuencias conocidas:
\begin{itemize}[leftmargin=1.2cm]
  \item retrasos recurrentes en la generación de reportes ejecutivos,
  \item errores de transcripción al consolidar datos de distintas fuentes,
  \item escasa visibilidad en tiempo real sobre los indicadores del negocio,
  \item dependencia operativa de personas específicas para tareas repetitivas.
\end{itemize}

La arquitectura propuesta apunta a reducir el tiempo de generación de reportes y mejorar la confiabilidad de los datos que llegan a la dirección.

% ─────────────────────────────────────────────────────────────
\section{Alcance y Cumplimiento de Parcial 1}

\subsection{Alcance}

Este documento cubre la definición de microservicios clave, el diseño del API Gateway, la justificación de los patrones aplicados, la estrategia de escalabilidad, los diagramas de arquitectura y BPMN, los modelos de datos, los requerimientos funcionales y no funcionales, y los lineamientos de observabilidad.

\newpage
%--------------------------------------------
\subsection{Matriz de cumplimiento}

\begin{longtable}{p{0.34\textwidth}p{0.30\textwidth}p{0.30\textwidth}}
\toprule
\textbf{Exigencia} & \textbf{Decisión} & \textbf{Evidencia en el proyecto} \\
\midrule
Definir microservicios clave & 4 servicios backend desacoplados & \texttt{docker-compose.yml}, secciones 5 y 6 \\
Diseñar API Gateway & Enrutamiento, auth request, balanceo y rate limiting & \texttt{gateway/nginx.conf} \\
Repository Pattern & Persistencia con Spring Data JPA & \texttt{authservice/.../repository}, 
\texttt{data-ingestion-service/} \\
Factory Method & Selección de conectores por sistema fuente & \texttt{ConnectorFactory} \\
Circuit Breaker & Tolerancia a fallos con Resilience4j & \texttt{HttpSourceConnector}, \texttt{application.yml} \\
Escalabilidad y desacoplamiento & Réplicas de auth e ingestion, contratos REST & \texttt{docker-compose.yml}, API por servicio \\
Documentar decisiones y justificar patrones & ADR técnico y análisis de riesgos & Secciones 7, 8 y 11 \\
\bottomrule
\end{longtable}


\newpage
% ─────────────────────────────────────────────────────────────
\section{Requerimientos}

\subsection{Requerimientos funcionales}
\begin{enumerate}[leftmargin=1.2cm]
  \item \textbf{RF-01:} Registrar y autenticar usuarios con tokens JWT.
  \item \textbf{RF-02:} Validar el token de acceso para proteger endpoints internos.
  \item \textbf{RF-03:} Ingerir datos desde sistemas fuente heterogéneos.
  \item \textbf{RF-04:} Persistir eventos y trazas de ejecución de sincronización.
  \item \textbf{RF-05:} Notificar al motor KPI luego de cada ingesta exitosa.
  \item \textbf{RF-06:} Centralizar el acceso mediante API Gateway.
  \item \textbf{RF-07:} Exponer el estado operativo de los servicios para monitoreo.
\end{enumerate}

\subsection{Requerimientos no funcionales}
\begin{enumerate}[leftmargin=1.2cm]
  \item \textbf{RNF-01 (Escalabilidad):} despliegue horizontal de servicios stateless.
  \item \textbf{RNF-02 (Disponibilidad):} tolerancia a fallos de fuentes externas mediante Circuit Breaker.
  \item \textbf{RNF-03 (Seguridad):} autenticación JWT, validación centralizada y filtros de tráfico.
  \item \textbf{RNF-04 (Observabilidad):} logs estructurados, correlación de request y health checks.
  \item \textbf{RNF-05 (Mantenibilidad):} separación por dominio y contratos API versionables.
  \item \textbf{RNF-06 (Performance):} límite de tasa en gateway y timeouts configurables por fuente.
\end{enumerate}

% ─────────────────────────────────────────────────────────────
\section{Descripción de Microservicios}

\begin{table}[H]
\centering
\begin{tabularx}{\textwidth}{p{0.18\textwidth}p{0.26\textwidth}p{0.24\textwidth}Y}
\toprule
\textbf{Servicio} & \textbf{Responsabilidad} & \textbf{Tecnología principal} & \textbf{Estado} \\
\midrule
API Gateway & Punto único de entrada, ruteo, balanceo, auth request, filtros WAF y rate limit & NGINX 1.27 & Implementado \\
Auth Service & Registro, login, refresh, validación JWT, perfil autenticado & Spring Boot, Spring Security, JPA, PostgreSQL & Implementado \\
Data Ingestion Service & Ingesta por fuente, persistencia de eventos, registro de corridas, notificación a KPI & Spring Boot, WebClient, Resilience4j, JPA, PostgreSQL & Implementado \\
KPI Engine & Recálculo de indicadores y publicación de métricas de negocio & API REST + almacenamiento KPI & Propuesto (siguiente etapa) \\
\bottomrule
\end{tabularx}
\caption{Microservicios definidos para la solución}
\end{table}
\newpage
\subsection{Especificación resumida por servicio}

\subsubsection*{API Gateway}
\begin{itemize}[leftmargin=1.2cm]
  \item \textbf{Entrada:} HTTP en puerto 8080.
  \item \textbf{Rutas clave:} \texttt{/api/auth/*}, \texttt{/api/ingestion/*}.
  \item \textbf{Funciones no funcionales:} balanceo \textit{least-conn}, \textit{auth\_request}, \textit{limit\_req}, detección de firmas maliciosas en URI/query/user-agent.
\end{itemize}

\subsubsection*{Auth Service}
\begin{itemize}[leftmargin=1.2cm]
  \item \textbf{Endpoints:} \texttt{/api/auth/register}, \texttt{/api/auth/login}, \texttt{/api/auth/refresh}, \texttt{/api/auth/validate}, \texttt{/api/users/me}.
  \item \textbf{Modelo de seguridad:} JWT con Spring Security stateless.
  \item \textbf{Persistencia:} usuarios y refresh tokens en PostgreSQL.
\end{itemize}

\subsubsection*{Data Ingestion Service}
\begin{itemize}[leftmargin=1.2cm]
  \item \textbf{Endpoint principal:} \texttt{/api/ingestion/sync/\{sourceSystem\}}.
  \item \textbf{Patrones implementados:} Factory Method para conectores y Circuit Breaker por fuente.
  \item \textbf{Persistencia:} eventos en \texttt{ingestion\_events} y trazas en \texttt{sync\_runs}.
\end{itemize}

\subsubsection*{KPI Engine (propuesto)}
\begin{itemize}[leftmargin=1.2cm]
  \item \textbf{Endpoint esperado:} \texttt{/api/kpi/recalculate}.
  \item \textbf{Contrato de entrada:} sistema fuente afectado y cantidad de registros procesados.
  \item \textbf{Salida esperada:} snapshot de indicadores actualizado.
\end{itemize}

% ─────────────────────────────────────────────────────────────
\section{Arquitectura General}

\begin{figure}[H]
\centering
\begin{tikzpicture}[
  font=\small,
  node distance=9mm and 14mm,
  % Estilos mejorados
  client/.style={
    draw=borderorange, line width=1.2pt,
    rounded corners=4pt, align=center,
    minimum height=14mm, minimum width=34mm,
    fill=boxorange, drop shadow
  },
  svc/.style={
    draw=borderblue, line width=1.2pt,
    rounded corners=4pt, align=center,
    minimum height=12mm, minimum width=34mm,
    fill=boxblue, drop shadow
  },
  propuesto/.style={
    draw=borderblue, line width=1pt, dashed,
    rounded corners=4pt, align=center,
    minimum height=12mm, minimum width=34mm,
    fill=blue!3
  },
  db/.style={
    draw=bordergreen, line width=1.2pt,
    cylinder, shape border rotate=90, aspect=0.3,
    minimum height=14mm, minimum width=18mm,
    align=center, fill=boxgreen, drop shadow
  },
  source/.style={
    draw=borderorange, line width=1pt,
    rounded corners=4pt, align=center,
    minimum height=11mm, minimum width=38mm,
    fill=boxorange!60
  },
  arr/.style={-{Latex[length=2.5mm]}, line width=1.2pt, color=arrowgray},
  notifyarr/.style={-{Latex[length=2.5mm]}, line width=1pt, dashed, color=borderblue},
  label/.style={font=\scriptsize\itshape, color=arrowgray}
]

% Nodos
\node[client]  (frontend) {Frontend Web\\Dashboard};
\node[svc, right=20mm of frontend] (gateway) {API Gateway\\{\footnotesize(NGINX 1.27)}};
\node[svc, above right=10mm and 18mm of gateway] (auth) {Auth Service\\{\footnotesize(Spring Boot)}};
\node[svc, right=20mm of gateway] (ing) {Data Ingestion\\Service};
\node[propuesto, below right=10mm and 18mm of gateway] (kpi) {KPI Engine\\{\footnotesize(propuesto)}};

\node[db, right=14mm of auth]  (authdb)   {Auth\\DB};
\node[db, right=14mm of ing]   (eventsdb) {Events\\DB};
\node[db, right=14mm of kpi]   (kpidb)    {KPI\\DB};

\node[source, below=16mm of ing] (sources) {Sistemas Fuente\\{\footnotesize POS · ERP · CRM · Finanzas}};

% Flechas
\draw[arr] (frontend) -- node[label, above]{HTTPS} (gateway);
\draw[arr] (gateway)  -- node[label, above left]{\texttt{/api/auth}} (auth);
\draw[arr] (gateway)  -- node[label, above]{\texttt{/api/ingestion}} (ing);
\draw[notifyarr] (ing) -- node[label, right]{notify} (kpi);
\draw[arr] (auth)   -- (authdb);
\draw[arr] (ing)    -- (eventsdb);
\draw[arr] (kpi)    -- (kpidb);
\draw[arr] (sources) -- (ing);

% Leyenda
\node[font=\scriptsize, anchor=south west] at (current bounding box.south west) 
  {\textcolor{borderblue}{$\longrightarrow$} flujo directo \quad \textcolor{borderblue}{$\dashrightarrow$} notificación asíncrona};

\end{tikzpicture}
\caption{Arquitectura de microservicios --- Plataforma de Monitoreo Inteligente}
\end{figure}

% ─────────────────────────────────────────────────────────────
\section{Orquestación con Docker Compose}

La orquestación local usa \texttt{docker-compose.yml} con dos réplicas de \texttt{authservice}, dos réplicas de \texttt{data-ingestion-service}, el API Gateway NGINX, dos bases PostgreSQL independientes (una por dominio) y volúmenes persistentes.

\begin{figure}[H]
\centering
\begin{tikzpicture}[
  font=\small,
  node distance=5mm and 10mm,
  gw/.style={
    draw=borderblue, line width=1.4pt, rounded corners=4pt,
    align=center, minimum width=32mm, minimum height=11mm,
    fill=boxblue, drop shadow
  },
  svc/.style={
    draw=borderblue, line width=1pt, rounded corners=3pt,
    align=center, minimum width=34mm, minimum height=10mm,
    fill=boxblue!60
  },
  db/.style={
    draw=bordergreen, line width=1pt,
    cylinder, shape border rotate=90, aspect=0.3,
    minimum height=13mm, minimum width=14mm,
    align=center, fill=boxgreen
  },
  grp/.style={
    draw=borderblue!50, line width=0.8pt, dashed,
    rounded corners=6pt, inner sep=5mm,
    fill=lanecolor
  },
  arr/.style={-{Latex[length=2mm]}, line width=1pt, color=arrowgray}
]

% Gateway
\node[gw] (gw) {api-gateway\\{\footnotesize :8080}};

% Auth replicas
\node[svc, right=22mm of gw, yshift=13mm]  (a1) {authservice-1\\{\footnotesize :8080}};
\node[svc, below=4mm of a1]                (a2) {authservice-2\\{\footnotesize :8080}};
\node[db,  right=12mm of a1, yshift=-3mm]  (dba) {postgres\\auth};

% Ingestion replicas
\node[svc, right=22mm of gw, yshift=-18mm] (i1) {ingestion-1\\{\footnotesize :8081}};
\node[svc, below=4mm of i1]                (i2) {ingestion-2\\{\footnotesize :8081}};
\node[db,  right=12mm of i1, yshift=-3mm]  (dbi) {postgres\\events};

% KPI (propuesto)
\node[svc, draw=borderblue!40, fill=blue!4, dashed, below=8mm of i2] (kpi) {kpi-engine\\{\footnotesize :8082 (target)}};

% Grupos
\begin{scope}[on background layer]
  \node[grp, fit=(a1)(a2)(dba), label={[font=\footnotesize\bfseries, color=laneheader]above:Dominio Autenticación}] {};
  \node[grp, fit=(i1)(i2)(dbi)(kpi), label={[font=\footnotesize\bfseries, color=laneheader]below:Dominio Ingesta y KPI}] {};
\end{scope}

% Flechas
\draw[arr] (gw) -- (a1);
\draw[arr] (gw) -- (a2);
\draw[arr] (gw) -- (i1);
\draw[arr] (gw) -- (i2);
\draw[arr] (a1) -- (dba);
\draw[arr] (a2) -- (dba);
\draw[arr] (i1) -- (dbi);
\draw[arr] (i2) -- (dbi);
\draw[arr, dashed] (i1) -- (kpi);
\draw[arr, dashed] (i2) -- (kpi);

\end{tikzpicture}
\caption{Orquestación de servicios con docker-compose y réplicas por dominio}
\end{figure}
\newpage
% ─────────────────────────────────────────────────────────────
\section{BPMN del Flujo End-to-End}

El diagrama muestra el ciclo completo de una sincronización: desde la recepción del request en el gateway hasta el recálculo de KPI y el ACK de vuelta.

\begin{figure}[H]
\centering
\begin{tikzpicture}[
  font=\scriptsize,
  scale=0.72,
  transform shape,
  node distance=3mm and 4mm,
  task/.style={
    draw=borderblue, line width=0.9pt,
    rounded corners=3pt,
    minimum width=26mm, minimum height=7mm,
    align=center, fill=boxblue
  },
  taskred/.style={
    draw=red!60, line width=0.9pt,
    rounded corners=3pt,
    minimum width=22mm, minimum height=7mm,
    align=center, fill=red!8
  },
  event/.style={
    draw=bordergreen, line width=1.2pt,
    circle, minimum size=6mm, fill=boxgreen
  },
  endevent/.style={
    draw=bordergreen, line width=2pt,
    circle, minimum size=6mm, fill=boxgreen!60
  },
  gw/.style={
    draw=borderorange, line width=1pt,
    diamond, aspect=2, align=center,
    minimum width=18mm, minimum height=8mm,
    fill=boxorange, inner sep=0.5pt
  },
  flow/.style={-{Latex[length=2mm]}, line width=0.9pt, color=arrowgray},
  dflow/.style={-{Latex[length=2mm]}, line width=0.8pt, dashed, color=borderblue!70},
  lanelabel/.style={font=\scriptsize\bfseries, color=laneheader, rotate=90, anchor=center}
]

% ── Fondo de lanes ──
\fill[lanecolor!80] (0,9) rectangle (24,12);
\fill[lanecolor!50] (0,6) rectangle (24,9);
\fill[lanecolor!80] (0,3) rectangle (24,6);
\fill[lanecolor!50] (0,0) rectangle (24,3);

% ── Bordes ──
\draw[line width=1pt, color=borderblue!40] (0,0) rectangle (24,12);
\foreach \y in {3,6,9} \draw[color=borderblue!30] (0,\y) -- (24,\y);
\draw[line width=1.5pt, color=borderblue!40] (1.5,0) -- (1.5,12);

% ── Etiquetas de lane ──
\node[lanelabel] at (0.75,10.5) {API Gateway};
\node[lanelabel] at (0.75,7.5)  {Auth Service};
\node[lanelabel] at (0.75,4.5)  {Data Ingestion};
\node[lanelabel] at (0.75,1.5)  {KPI Engine};

% ── Lane: API Gateway ──
\node[event]  (s)    at (2.4,10.5) {};
\node[task]   (g1)   at (5.5,10.5) {Recibe request\\[-1pt]\texttt{/api/ingestion/sync}};
\node[task]   (g2)   at (10.2,10.5){Envía \texttt{auth\_request}\\[-1pt]a \texttt{/auth/validate}};
\node[gw]     (g3)   at (14.8,10.5){Token\\válido?};
\node[task]   (g4)   at (19.5,10.5){Proxy a\\Data Ingestion};
\node[endevent](e)   at (22.8,10.5){};
\node[taskred] (deny) at (19.5,9.2) {401 Unauthorized};

% ── Lane: Auth Service ──
\node[task] (a1) at (10.2,7.5) {Valida JWT};
\node[task] (a2) at (14.8,7.5) {Responde 200/401};

% ── Lane: Data Ingestion ──
\node[task] (i1) at (19.5,4.5) {Selecciona conector\\[-1pt](Factory Method)};
\node[gw]   (i2) at (14.8,4.5) {CB\\abierto?};
\node[task] (i3) at (10.2,4.5) {Extrae y transforma\\[-1pt]eventos};
\node[task] (i4) at (6.2,4.5)  {Persiste en\\events\_db};
\node[task] (i5) at (2.8,4.5)  {Notifica\\recálculo KPI};

% ── Lane: KPI Engine ──
\node[task] (k1) at (2.8,1.5)  {Recalcula KPI};
\node[task] (k2) at (6.2,1.5)  {Guarda snapshot};
\node[task] (k3) at (10.2,1.5) {ACK al emisor};

% ── Flujos ──
\draw[flow] (s) -- (g1);
\draw[flow] (g1) -- (g2);
\draw[flow] (g2) -- (g3);
\draw[flow] (g3) -- node[above, font=\tiny]{Sí} (g4);
\draw[flow] (g3) -- node[right, font=\tiny]{No} (deny);
\draw[flow] (g4) -- (e);

\draw[dflow] (g2) -- (a1);
\draw[dflow] (a1) -- (a2);
\draw[dflow] (a2) -- (g3);

\draw[dflow] (g4) -- (i1);
\draw[flow]  (i1) -- (i2);
\draw[flow]  (i2) -- node[above, font=\tiny]{No} (i3);
\draw[flow]  (i3) -- (i4);
\draw[flow]  (i4) -- (i5);
\draw[flow]  (i2) |- node[left, font=\tiny]{Sí (fallback)} (i4);

\draw[dflow] (i5) -- (k1);
\draw[flow]  (k1) -- (k2);
\draw[flow]  (k2) -- (k3);
\draw[dflow] (k3) -| (g4);

% Leyenda evento inicio/fin
\node[font=\tiny, anchor=south] at (2.4,11.2) {inicio};
\node[font=\tiny, anchor=south] at (22.8,11.2) {fin};

\end{tikzpicture}
\caption{BPMN simplificado del flujo end-to-end entre los 4 servicios. Línea continua = flujo de control interno; línea punteada = llamada entre servicios.}
\end{figure}

\newpage

% ─────────────────────────────────────────────────────────────
\section{Modelos de Datos por Servicio}

\subsection{Auth Service}

\begin{figure}[H]
\centering
\begin{tikzpicture}[
  font=\small,
  ent/.style={
    draw=borderblue, line width=1.2pt,
    rounded corners=4pt,
    minimum width=62mm, align=left,
    fill=boxblue, drop shadow, inner sep=4pt
  },
  pk/.style={font=\small\bfseries},
  fk/.style={font=\small\itshape},
  rel/.style={-{Latex[length=2.5mm]}, line width=1.2pt, color=arrowgray}
]

\node[ent] (users) {
  \textbf{users}\\[2pt]
  \texttt{\textbf{PK} id: bigint}\\
  \texttt{email: varchar(320) UNIQUE}\\
  \texttt{password\_hash: varchar}\\
  \texttt{role: enum(USER, ADMIN)}\\
  \texttt{created\_at, updated\_at: instant}
};

\node[ent, right=30mm of users] (rt) {
  \textbf{refresh\_tokens}\\[2pt]
  \texttt{\textbf{PK} id: bigint}\\
  \texttt{token: varchar(512) UNIQUE}\\
  \texttt{\textit{FK} user\_id $\to$ users.id}\\
  \texttt{expires\_at: instant}\\
  \texttt{revoked: boolean}\\
  \texttt{created\_at: instant}
};

\draw[rel] (users.east) -- node[above, font=\scriptsize]{1 $\cdots$ N} (rt.west);

\end{tikzpicture}
\caption{Modelo de datos --- Auth Service}
\end{figure}

\subsection{Data Ingestion Service}

\begin{figure}[H]
\centering
\begin{tikzpicture}[
  font=\small,
  ent/.style={
    draw=bordergreen, line width=1.2pt,
    rounded corners=4pt,
    minimum width=64mm, align=left,
    fill=boxgreen, drop shadow, inner sep=4pt
  }
]

\node[ent] (ev) {
  \textbf{ingestion\_events}\\[2pt]
  \texttt{\textbf{PK} id: bigint}\\
  \texttt{source\_system: varchar(80)}\\
  \texttt{ingested\_at: timestamptz}\\
  \texttt{payload: jsonb}
};

\node[ent, right=22mm of ev] (sr) {
  \textbf{sync\_runs}\\[2pt]
  \texttt{\textbf{PK} id: bigint}\\
  \texttt{source\_system: varchar(80)}\\
  \texttt{records\_processed: int}\\
  \texttt{started\_at, completed\_at: timestamptz}\\
  \texttt{status: varchar(30)}
};

\end{tikzpicture}
\caption{Modelo de datos --- Data Ingestion Service}
\end{figure}

\subsection{KPI Engine (propuesto)}

\begin{figure}[H]
\centering
\begin{tikzpicture}[
  font=\small,
  ent/.style={
    draw=borderorange, line width=1.2pt,
    rounded corners=4pt,
    minimum width=66mm, align=left,
    fill=boxorange, drop shadow, inner sep=4pt
  },
  rel/.style={-{Latex[length=2.5mm]}, line width=1.2pt, color=arrowgray}
]

\node[ent] (kd) {
  \textbf{kpi\_definitions}\\[2pt]
  \texttt{\textbf{PK} id: bigint}\\
  \texttt{code: varchar(40) UNIQUE}\\
  \texttt{name: varchar(120)}\\
  \texttt{formula: text}\\
  \texttt{frequency: varchar(20)}
};

\node[ent, right=26mm of kd] (ks) {
  \textbf{kpi\_snapshots}\\[2pt]
  \texttt{\textbf{PK} id: bigint}\\
  \texttt{\textit{FK} kpi\_id $\to$ kpi\_definitions.id}\\
  \texttt{period\_start, period\_end: timestamptz}\\
  \texttt{value: numeric}\\
  \texttt{computed\_at: timestamptz}
};

\draw[rel] (kd.east) -- node[above, font=\scriptsize]{1 $\cdots$ N} (ks.west);

\end{tikzpicture}
\caption{Modelo de datos propuesto --- KPI Engine}
\end{figure}

\subsection{Modelo lógico del API Gateway}

\begin{table}[H]
\centering
\begin{tabularx}{\textwidth}{p{0.22\textwidth}Y Y}
\toprule
\textbf{Entidad lógica} & \textbf{Atributos clave} & \textbf{Descripción} \\
\midrule
RouteRule & path\_prefix, auth\_required, rate\_limit & Reglas de ruteo y seguridad por prefijo URI \\
UpstreamGroup & strategy, servers, fail\_timeout, retries & Cluster destino y política de balanceo/failover \\
SecurityFilter & uri\_signature, query\_signature, user\_agent\_signature & Bloqueo básico de tráfico malicioso \\
TraceHeader & correlation\_id, forwarded\_for & Propagación de trazabilidad entre servicios \\
\bottomrule
\end{tabularx}
\caption{Modelo lógico no persistente del API Gateway}
\end{table}

% ─────────────────────────────────────────────────────────────
\section{Patrones de Diseño Aplicados}

\subsection{Repository Pattern}

Spring Data JPA encapsula toda la lógica de persistencia para usuarios, tokens, eventos y corridas de sincronización. El beneficio directo es que la lógica de negocio no sabe nada de SQL: se puede cambiar el motor de base de datos sin tocar los servicios. También hace que las pruebas unitarias sean directas, ya que los repositorios se pueden sustituir por mocks.

\subsection{Factory Method}

\texttt{ConnectorFactory} construye el conector correcto según el sistema fuente indicado en el request. Agregar un nuevo sistema de origen (por ejemplo, un nuevo proveedor de e-commerce) implica escribir una clase de conector y registrarla en el factory; la lógica de ingesta principal no cambia.

\subsection{Circuit Breaker}

Resilience4j protege cada conector de forma independiente. Si un sistema fuente empieza a fallar o a responder lento, el circuito se abre y el servicio devuelve un fallback controlado en lugar de acumular hilos bloqueados. Cuando la fuente se recupera, Resilience4j hace la transición al estado semiabierto para verificar antes de reanudar el tráfico completo.

\subsection{Patrones complementarios}

\begin{itemize}[leftmargin=1.2cm]
  \item \textbf{API Gateway Pattern:} concentra seguridad, ruteo y políticas transversales en un único punto de entrada.
  \item \textbf{DTO + Validation:} contratos explícitos con validación temprana, antes de que el dato entre al dominio.
  \item \textbf{Stateless Services:} al no guardar estado en memoria, cada réplica es equivalente, lo que hace que el escalado horizontal sea trivial.
\end{itemize}

% ─────────────────────────────────────────────────────────────
\section{Decisiones Técnicas}

\begin{longtable}{p{0.20\textwidth}p{0.25\textwidth}p{0.22\textwidth}p{0.27\textwidth}}
\toprule
\textbf{Decisión} & \textbf{Alternativas} & \textbf{Selección} & \textbf{Justificación} \\
\midrule
Arquitectura backend & Monolito modular vs microservicios & Microservicios & Independencia de despliegue y escalado por dominio \\
Gateway & Spring Cloud Gateway vs NGINX & NGINX & Menor huella operativa, excelente rendimiento en proxy y balanceo \\
Persistencia & DB única vs DB por dominio & DB por dominio (auth/events) & Aislamiento funcional y menor acoplamiento entre esquemas \\
Resiliencia & Reintentos simples vs CB+timeouts & Circuit Breaker + timeout & Reduce el impacto de fuentes inestables y evita saturación del pool \\
Autenticación & Sesión stateful vs JWT stateless & JWT stateless & Escalabilidad horizontal sin dependencia de memoria compartida \\
Contratos & Integración directa por BD vs API REST & API REST & Encapsulación de dominio y control de cambios por versión \\
\bottomrule
\end{longtable}

% ─────────────────────────────────────────────────────────────
\section{Escalabilidad, Desacoplamiento y Seguridad}

\subsection{Escalabilidad}

Auth e Ingestion corren con dos réplicas desde el día uno. El gateway distribuye tráfico con \textit{least-conn}, por lo que agregar una tercera réplica no requiere cambios de configuración. Los servicios no guardan estado en memoria, así que cada instancia nueva es funcional de inmediato.

\subsection{Desacoplamiento}

Los servicios solo se hablan por contratos REST; ninguno accede directamente a la base de datos del otro. La configuración de fuentes y endpoints es externa. Esto permite evolucionar un servicio sin coordinar cambios en los demás, salvo que se modifique el contrato de la API.

\subsection{Seguridad}

El JWT se valida en el Auth Service y el gateway refuerza ese control con \texttt{auth\_request} antes de enrutar cualquier request interno. A nivel de red, el gateway aplica rate limiting por IP, reglas básicas de firma para bloquear patrones maliciosos en URI/query/user-agent, y CORS controlado. Los servicios internos no están expuestos directamente.

% ─────────────────────────────────────────────────────────────
\section{Observabilidad}

\begin{enumerate}[leftmargin=1.2cm]
  \item \textbf{Correlation ID:} se genera en el gateway y se propaga como header a todos los servicios, lo que permite reconstruir el recorrido de un request en los logs.
  \item \textbf{Logging operativo:} NGINX registra acceso y errores; los servicios Spring usan logs estructurados.
  \item \textbf{Health checks:} endpoint de salud en Ingestion y healthchecks Docker para PostgreSQL.
  \item \textbf{Spring Actuator:} dependencia incorporada para métricas y readiness, con exposición configurable por perfil.
  \item \textbf{Trazabilidad de negocio:} cada sincronización queda registrada en \texttt{sync\_runs} con estado y tiempos, lo que permite auditar qué procesó cada corrida.
\end{enumerate}

% ─────────────────────────────────────────────────────────────
\section{Riesgos Técnicos y Mitigaciones}

\begin{table}[H]
\centering
\begin{tabularx}{\textwidth}{p{0.26\textwidth}p{0.20\textwidth}Y}
\toprule
\textbf{Riesgo} & \textbf{Impacto} & \textbf{Mitigación propuesta} \\
\midrule
KPI Engine no desplegado aún & Medio-Alto & Mantener el contrato REST definido y usar un mock para pruebas de integración hasta la siguiente etapa \\
Fuentes externas inestables & Alto & Circuit Breaker con timeout por fuente y fallback controlado \\
Divergencia de versiones Java/Spring entre servicios & Medio & Normalizar baseline de runtime y BOM en un sprint técnico dedicado \\
Sobrecarga de tráfico en el gateway & Medio & Rate limit, escalado horizontal y monitoreo de latencia p95/p99 \\
\bottomrule
\end{tabularx}
\caption{Principales riesgos de arquitectura en la etapa Parcial 1}
\end{table}

% ─────────────────────────────────────────────────────────────
\section{Conclusión}

La arquitectura entregada cubre lo requerido en Parcial 1: cuatro microservicios con responsabilidades bien delimitadas, un API Gateway funcional, los tres patrones solicitados implementados y evidenciados en código, y una base que puede crecer sin grandes refactorizaciones.

Hay un punto que vale la pena anticipar: el KPI Engine es el componente que más va a determinar si la plataforma cumple su promesa de visibilidad ejecutiva. Tiene un contrato definido y puede probarse con un mock hoy, pero su implementación real debería ser la prioridad de la siguiente etapa. Todo lo demás ya está en condiciones de producción.

% ─────────────────────────────────────────────────────────────
\section*{Anexo A: Extracto de orquestación}

\begin{lstlisting}[style=yamlstyle,caption={Extracto representativo de docker-compose.yml}]
services:
  authservice-1:
    build: ./authservice
    expose: ["8080"]

  authservice-2:
    build: ./authservice
    expose: ["8080"]

  data-ingestion-service-1:
    build: ./data-ingestion-service
    ports: ["8081:8081"]

  data-ingestion-service-2:
    build: ./data-ingestion-service
    expose: ["8081"]

  api-gateway:
    image: nginx:1.27-alpine
    ports: ["8080:8080"]
\end{lstlisting}

\end{document}